\subsection{The actual Case II modular ideal}
\label{subsec:sr-case-II-actual-pm}

We apply the common module argument in rank four only at its
stable curves. No rank-three component graph is used.

The group of the common space in this subsection is \(G_F\),
where \(F\) is the tower's totally real field and \(L/F\) is
its quadratic CM extension; it is not the original rational group.
Fix the actual tame Case II tower, its two links \(K,K'\) with
the same scalar character \(\kappa\), and their distinct neat
places \(v_0,v'_0\).
Let \(X=(A,B,P)\) be its common rank-four space.
The scalar-twisted restriction of \(P\) to \(G_L\), with its
prescribed ordinary characters, gives by the universal determinant
property the maps
\[
 \phi:\mathcal B_K\longrightarrow B,\qquad
 \phi':\mathcal B_{K'}\longrightarrow B.
\]
Their formulas are the characteristic polynomial of the
scalar-twisted restriction and its four ordered characters;
the injective coefficient map \(B\hookrightarrow A\) makes
the determinant and ordinary identities hold already over \(B\).
The compact integral interpolation argument used for (H9)
applies equally here and does not divide by \(4!\).
Let \(J_K,J_{K'}\) be the actual Hecke kernels of
Proposition~\ref{prop:sr-case-two-hecke-quotient}. Define
\[
                  J=\phi(J_K)B+\phi'(J_{K'})B.             \tag{R4M1}
\]
This is one ideal in the actual common space.

\paragraph{The rank-four stable-curve support input.}
For either link repeat (CA6)--(CA13) with \(n=4\), using
the coefficient system and classical realization
(H1)--(H11) already proved in
Section~\ref{subsec:sr-case-two-hecke}.
The field assumptions are satisfied: the CM extension is
finite-unramified, every prescribed ramified place and every
\(\ell\)-place splits, and \(4\mid4[F:\mathbb Q]\).
Thorne's Theorem 6.5 therefore includes \(\ell=3\).
The ordinary function modules, exact classical control and
auxiliary projectivity are rank-independent.
The same local double-coset maps and the finite torus weight
calculation give the uniform exponent \(4!=24\).
All auxiliary levels have the same tame coefficient lattice
and the same base Hecke kernel.

For an actual characteristic-\(\ell\) curve in \(X\) with
irreducible generic representation and distinct ordered local
characters on \(U_v\), the link construction gives its
symmetric polarized determinant point.
Proposition~\ref{prop:sr-four-curve-genericity} supplies the
precise noncommutative image, cyclotomic restriction and
regular-coset conditions required for this point's
Taylor--Wiles prime choices.
The actual tower's root separators and distinguished scalar
Frobenius give its other local conditions, including \(N4\);
these are the witnesses checked in the tower construction.
The first-Selmer and prime-choice proofs are rank-independent
under exactly these hypotheses.

Apply the one common patched module and the same-completion
calculation of Section~\ref{subsec:sr-same-completion-galout},
now with
\[
                   D_0=4d+4q_{\rm TW}.
\]
The point is stable, so its completed Cayley--Hamilton algebra
is the one-block algebra and its odd module is zero.
Only the ordinary presentation-prime lemma is needed.
It makes the stable completed support component contain every
prime below the curve point.
Transport to \(B\) proves that every prime beneath the curve
kernel contains \(\phi(J_K)B\); the second link proves the
same assertion for \(\phi'(J_{K'})B\).
Thus \(J\) has exactly the stable-curve local-zero property
required in Proposition~\ref{prop:sr-four-core}.
This uses the rank-four common module and its stable-point
presentation, not the rank-two/rank-three propagation theorem.

\paragraph{Finite weight support.}
We verify \eqref{eq:sr-pm-finite} without assuming that either
map \(\phi\) is surjective.
Let \(\mathfrak q\subset B\) contain \(J\), the coefficient
uniformizer, and every \(\psi_{vi}(u)-1\) for inertia \(u\)
above \(\ell\). Put \(R=B/\mathfrak q\).
The scalar link character is fixed.
Consequently all inertia weight characters of the determinant
pair over \(R\) are constants.
Universality of the local weight ring makes its entire image
constant, by the root-of-unity factorization used in (SF4).

The Hecke quotient is finite over that weight ring.
Every image under \(\phi\) of a determinant coefficient or
ordinary character is therefore integral over the finite
coefficient field.
The fixed scalar twist is a unit constant here and can be
removed. Thus the characteristic coefficients on \(G_L\)
are integral over this field.
For \(g\in G_F\), \(g^2\in G_L\).
More generally, for any finite-index restriction one may use a
positive power lying in the restricting subgroup; the argument
below uses no bound on this exponent.
Each eigenvalue of \(P(g)\) is integral once its square is
integral; hence every characteristic coefficient of \(P(g)\)
is integral.
At every ordinary local group the CM place splits, so its
local characters are already obtained by restriction from
the link pair and are integral as well.

In the local domain \(R\), an element algebraic over its
finite residue field equals its specified residue constant.
Indeed a nonzero algebraic element is a unit, whereas the
difference from its residue lies in the maximal ideal;
equivalently use the separable finite-field minimal polynomial.
The common space's characteristic-coefficient and
ordinary-character generation property therefore makes
\(B\to R\) equal to its residue map.
Thus \(\mathfrak q=\mathfrak m_B\).
Noetherianity now gives \eqref{eq:sr-pm-finite} by taking a
power of the maximal ideal.

\paragraph{The other three modular-ideal properties.}
Proposition~\ref{prop:sr-seed-local-zero}, applied to each
link and transported through \(\phi,\phi'\), says that \(J\)
is zero locally at the actual induced seed.
This is the seed property, stronger than merely having an
arithmetic point kill \(J\).
At a prescribed characteristic-zero weight, killing \(J\)
gives actual classical ordinary eigensystems through the two
Hecke quotients.
The strict spectral inequality in
Lemma~\ref{lem:sr-case-two-speh-exclusion} gives their purity
condition on the original representation.
If the restriction to \(G_L\) is irreducible, the classical
base change is cuspidal.
The exact descent and full local compatibility proved in
Proposition~\ref{prop:sr-all-quadratic-automorphy} and
Proposition~\ref{prop:sr-actual-wit} apply to it.
The two different neat places remove the one omitted local
place in either individual link comparison.
The common scalar \(\kappa\) is essential: it makes the
two resulting systems refer to the same rational
representation before gluing.

We have constructed all five fields of the actual
\srcid{PotentialModularityTower}: finite weight support,
stable-curve local zero, seed local zero, purity, and
automorphy at the specified irreducible arithmetic points.
The frozen-source assembly is \srcid{PmTwo}; the stable-curve
step is \srcid{ModularDeformationQuotientRegularAtPoint}.
This conclusion concerns an actual rank-four common space and
the two actual links. Construction of that space and of the
residual paths connecting its points to the seed is a separate
rank-four step; it does not follow from this modular-ideal
construction alone.
